% GENERATED FILE. Edit canonical lessons, then run npm run generate:books.
% canonical-source-hash: fnv1a64:6d6363deabb2b602
% canonical-lessons: MW-C170-dhono, MW-C170-nichovno, MW-C170-pasarno, MW-C170-sukavno, MW-C170-sametno

\chapter{Wash Day}
\label{ch:mw-wash-day}
\hlchaptermodality{170}

\begin{chapteropening}
\textbf{By the end of this chapter:} \emph{I can say how I wash, wring and dry the clothes.}

Complete the last lesson of chapter 170: I can say how I wash, wring and dry the clothes.
\end{chapteropening}

\section[\texorpdfstring{\mw{धोणो} (dhoṇo)}{धोणो (dhoṇo)}]{\mw{धोणो} (dhoṇo) — to wash}

\label{lesson:MW-C170-dhono}



\begin{quote}
Before the new one: say the Marwadi for to fry, then the Marwadi for to taste.
\end{quote}

\subsection*{You'll want to know: \mw{धोणो}}
\textbf{\mw{धोणो}} — \emph{dhoṇo} — \textquotedblleft{}to wash\textquotedblright{}.

\begin{quote}
\mw{हूं} \mw{रोज} \mw{कपड़ा} \mw{धोऊं} \mw{हूं।}
\end{quote}

\begin{quote}
\emph{hū̃ roj kapṛā dhoū̃ hū̃.}
\end{quote}

\begin{quote}
I wash clothes every day.
\end{quote}

Past (I …): \mw{म्हैं} \mw{कपड़ा} \mw{धोया}, \emph{mhaĩ kapṛā dhoyā}.

\begin{cousinweb}
\textbf{\mw{धोणो}} (\emph{dhoṇo}) is inherited from a reconstructed \emph{*dhauvati}, \textquotedblleft{}washes\textquotedblright{}; CDIAL lists Marwari \emph{dhoṇo} under it (CDIAL 6886).
\end{cousinweb}

\begin{grammarlens}[title={What you've built}]
One more word for saying what you do.
\end{grammarlens}

\subsection*{Your turn}
\begin{itemize}
\raggedright
  \item \emph{Say it:} \emph{dhoṇo}
  \item \emph{Say it:} \emph{dhoṇo}, once more
  \item \emph{Say it:} \emph{hū̃ roj kapṛā dhoū̃ hū̃.} — I wash clothes every day.
  \item \emph{Recall:} say the Marwadi for to fry, then the Marwadi for to taste, then say \emph{dhoṇo} again
\end{itemize}

\begin{tcolorbox}[breakable,colback=teal!4,colframe=teal!35!black,title={Before you move on}]
What does \mw{धोणो} mean? (\textquotedblleft{}To wash\textquotedblright{}.) Say it once more.
\end{tcolorbox}
\section[\texorpdfstring{\mw{निचोवणो} (nichovṇo)}{निचोवणो (nichovṇo)}]{\mw{निचोवणो} (nichovṇo) — to wring out}

\label{lesson:MW-C170-nichovno}



\begin{quote}
Before the new one: say the Marwadi for to taste, then the Marwadi for to wash.
\end{quote}

\subsection*{You'll want to know: \mw{निचोवणो}}
\textbf{\mw{निचोवणो}} — \emph{nichovṇo} — \textquotedblleft{}to wring out\textquotedblright{}.

\begin{quote}
\mw{मां} \mw{कपड़ा} \mw{निचोवै} \mw{है।}
\end{quote}

\begin{quote}
\emph{mā̃ kapṛā nichovai hai.}
\end{quote}

\begin{quote}
Mother wrings out the clothes.
\end{quote}

Past (I …): \mw{म्हैं} \mw{कपड़ा} \mw{निचोया}, \emph{mhaĩ kapṛā nichoyā}.

\begin{cousinweb}
\textbf{\mw{निचोवणो}} (\emph{nichovṇo}) comes from a reconstructed \emph{*niścyotayati}, \textquotedblleft{}causes to trickle out\textquotedblright{}; CDIAL lists Hindi \emph{niconā}, \textquotedblleft{}to wring out\textquotedblright{}, and Gujarati \emph{nicovvũ} (CDIAL 7451). CDIAL gives no Marwari form in this entry.
\end{cousinweb}

\begin{grammarlens}[title={What you've built}]
One more word for saying what you do.
\end{grammarlens}

\subsection*{Your turn}
\begin{itemize}
\raggedright
  \item \emph{Say it:} \emph{nichovṇo}
  \item \emph{Say it:} \emph{nichovṇo}, once more
  \item \emph{Say it:} \emph{mā̃ kapṛā nichovai hai.} — Mother wrings out the clothes.
  \item \emph{Recall:} say the Marwadi for to taste, then the Marwadi for to wash, then say \emph{nichovṇo} again
\end{itemize}

\begin{tcolorbox}[breakable,colback=teal!4,colframe=teal!35!black,title={Before you move on}]
What does \mw{निचोवणो} mean? (\textquotedblleft{}To wring out\textquotedblright{}.) Say it once more.
\end{tcolorbox}
\section[\texorpdfstring{\mw{पसारणो} (pasārṇo)}{पसारणो (pasārṇo)}]{\mw{पसारणो} (pasārṇo) — to spread out to dry; to stretch out}

\label{lesson:MW-C170-pasarno}



\begin{quote}
Before the new one: say the Marwadi for to wash, then the Marwadi for to wring out.
\end{quote}

\subsection*{You'll want to know: \mw{पसारणो}}
\textbf{\mw{पसारणो}} — \emph{pasārṇo} — \textquotedblleft{}to spread out to dry; to stretch out\textquotedblright{}.

\begin{quote}
\mw{हूं} \mw{छात} \mw{माथै} \mw{कपड़ा} \mw{पसारूं} \mw{हूं।}
\end{quote}

\begin{quote}
\emph{hū̃ chhāt māthai kapṛā pasārū̃ hū̃.}
\end{quote}

\begin{quote}
I spread the clothes out on the roof.
\end{quote}

Past (I …): \mw{म्हैं} \mw{कपड़ा} \mw{पसार्या}, \emph{mhaĩ kapṛā pasāryā}.

\begin{cousinweb}
\textbf{\mw{पसारणो}} (\emph{pasārṇo}) is inherited from Sanskrit \emph{prasārayati}, \textquotedblleft{}stretches out, extends\textquotedblright{}; CDIAL lists Marwari \emph{pasārṇo}, \textquotedblleft{}to spread out\textquotedblright{} (CDIAL 8838).
\end{cousinweb}

\begin{grammarlens}[title={What you've built}]
One more word for saying what you do.
\end{grammarlens}

\subsection*{Your turn}
\begin{itemize}
\raggedright
  \item \emph{Say it:} \emph{pasārṇo}
  \item \emph{Say it:} \emph{pasārṇo}, once more
  \item \emph{Say it:} \emph{hū̃ chhāt māthai kapṛā pasārū̃ hū̃.} — I spread the clothes out on the roof.
  \item \emph{Recall:} say the Marwadi for to wash, then the Marwadi for to wring out, then say \emph{pasārṇo} again
\end{itemize}

\begin{tcolorbox}[breakable,colback=teal!4,colframe=teal!35!black,title={Before you move on}]
What does \mw{पसारणो} mean? (\textquotedblleft{}To spread out to dry; to stretch out\textquotedblright{}.) Say it once more.
\end{tcolorbox}
\section[\texorpdfstring{\mw{सुकावणो} (sukāvṇo)}{सुकावणो (sukāvṇo)}]{\mw{सुकावणो} (sukāvṇo) — to dry (clothes)}

\label{lesson:MW-C170-sukavno}



\begin{quote}
Before the new one: say the Marwadi for to wring out, then the Marwadi for to spread out to dry.
\end{quote}

\subsection*{You'll want to know: \mw{सुकावणो}}
\textbf{\mw{सुकावणो}} — \emph{sukāvṇo} — \textquotedblleft{}to dry (clothes)\textquotedblright{}.

\begin{quote}
\mw{हूं} \mw{धूप} \mw{मांय} \mw{कपड़ा} \mw{सुकाऊं} \mw{हूं।}
\end{quote}

\begin{quote}
\emph{hū̃ dhūp mā̃y kapṛā sukāū̃ hū̃.}
\end{quote}

\begin{quote}
I dry the clothes in the sun.
\end{quote}

Past (I …): \mw{म्हैं} \mw{कपड़ा} \mw{सुकाया}, \emph{mhaĩ kapṛā sukāyā}.

\begin{cousinweb}
\textbf{\mw{सुकावणो}} (\emph{sukāvṇo}) is built on a reconstructed \emph{*śuṣkati}, \textquotedblleft{}becomes dry\textquotedblright{}; CDIAL lists Old Marwari \emph{sūkaï}, \textquotedblleft{}dries up\textquotedblright{}, and a Gujarati verb \emph{sukāvũ}, \textquotedblleft{}to dry\textquotedblright{} (CDIAL 12552). The adjective \textbf{\mw{सूको}}, dry, is its relative.
\end{cousinweb}

\begin{grammarlens}[title={What you've built}]
One more word for saying what you do.
\end{grammarlens}

\subsection*{Your turn}
\begin{itemize}
\raggedright
  \item \emph{Say it:} \emph{sukāvṇo}
  \item \emph{Say it:} \emph{sukāvṇo}, once more
  \item \emph{Say it:} \emph{hū̃ dhūp mā̃y kapṛā sukāū̃ hū̃.} — I dry the clothes in the sun.
  \item \emph{Recall:} say the Marwadi for to wring out, then the Marwadi for to spread out to dry, then say \emph{sukāvṇo} again
\end{itemize}

\begin{tcolorbox}[breakable,colback=teal!4,colframe=teal!35!black,title={Before you move on}]
What does \mw{सुकावणो} mean? (\textquotedblleft{}To dry (clothes)\textquotedblright{}.) Say it once more.
\end{tcolorbox}
\section[\texorpdfstring{\mw{समेटणो} (sameṭṇo)}{समेटणो (sameṭṇo)}]{\mw{समेटणो} (sameṭṇo) — to gather up; to fold away}

\label{lesson:MW-C170-sametno}



\begin{quote}
Before the new one: say the Marwadi for to spread out to dry, then the Marwadi for to dry (clothes).
\end{quote}

\subsection*{You'll want to know: \mw{समेटणो}}
\textbf{\mw{समेटणो}} — \emph{sameṭṇo} — \textquotedblleft{}to gather up; to fold away\textquotedblright{}.

\begin{quote}
\mw{आथण} \mw{मां} \mw{कपड़ा} \mw{समेटै} \mw{है।}
\end{quote}

\begin{quote}
\emph{āthaṇ mā̃ kapṛā sameṭai hai.}
\end{quote}

\begin{quote}
In the evening my mother gathers up the clothes.
\end{quote}

Past (I …): \mw{म्हैं} \mw{कपड़ा} \mw{समेट्या}, \emph{mhaĩ kapṛā sameṭyā}.

\begin{cousinweb}
\textbf{\mw{समेटणो}} (\emph{sameṭṇo}) is inherited from Sanskrit \emph{samveṣṭayati}, \textquotedblleft{}wraps up\textquotedblright{}; CDIAL lists Hindi \emph{sameṭnā}, \textquotedblleft{}to roll up\textquotedblright{}, and Gujarati \emph{sameṭvũ} (CDIAL 13026). CDIAL gives no Marwari form in this entry.
\end{cousinweb}

\begin{grammarlens}[title={What you've built}]
One more word for saying what you do.
\end{grammarlens}

\subsection*{Your turn}
\begin{itemize}
\raggedright
  \item \emph{Say it:} \emph{sameṭṇo}
  \item \emph{Say it:} \emph{sameṭṇo}, once more
  \item \emph{Say it:} \emph{āthaṇ mā̃ kapṛā sameṭai hai.} — In the evening my mother gathers up the clothes.
  \item \emph{Recall:} say the Marwadi for to spread out to dry, then the Marwadi for to dry (clothes), then say \emph{sameṭṇo} again
\end{itemize}

\begin{tcolorbox}[breakable,colback=teal!4,colframe=teal!35!black,title={Before you move on}]
What does \mw{समेटणो} mean? (\textquotedblleft{}To gather up; to fold away\textquotedblright{}.) Say it once more.
\end{tcolorbox}
